% PanAfriCon AI 2026 -- Author Rebuttal
% Submission 151: Beyond Detection Parity: Localization Parity in Pedestrian Detection
\documentclass[11pt]{article}

\usepackage[margin=1in]{geometry}
\usepackage[T1]{fontenc}
\usepackage{times}
\usepackage{enumitem}
\usepackage{xcolor}
\usepackage{hyperref}

\hypersetup{colorlinks=true,linkcolor=blue,urlcolor=blue,citecolor=blue}

\title{\textbf{Author Rebuttal}\\[0.4em]
\large PanAfriCon AI 2026 --- Submission 151\\[0.2em]
\normalsize Beyond Detection Parity: Localization Parity in Pedestrian Detection}
\author{Abel Adissu \quad Anwar Gashaw Yimam \quad Dawit Kindea \quad Beakal Gizachew \quad Elefelious Getachew\\[0.3em]
\textit{Addis Ababa University}\\[0.2em]
\small\texttt{abel.adissu-ug@aau.edu.et},
\texttt{anwar.gashaw-ug@aau.edu.et},\\
\texttt{dawit.kindea-ug@aau.edu.et},
\texttt{beakal.gizachew@aau.edu.et},\\
\texttt{elefelious.getachew@aau.edu.et}}
\date{}

\begin{document}
\maketitle
\thispagestyle{empty}

\noindent We thank the Program Committee and both reviewers for their careful
reading and constructive feedback. We are encouraged by the accept / weak-accept
recommendations and by Reviewer~1's recognition of Localization Parity as a
meaningful conceptual contribution. Below we address each comment directly.
All requested revisions have been incorporated into the camera-ready manuscript
(Springer LNCS / \texttt{llncs} format).

\vspace{0.6em}
\noindent\textbf{Summary of manuscript changes.}
\begin{itemize}[leftmargin=1.2em,itemsep=2pt]
\item Reformatted the full paper to the official Springer LNCS LaTeX template.
\item Named all datasets, detector families, and demographic axes explicitly in
      the abstract, introduction, and experimental setup.
\item Added a formal subsection relating Localization Parity to existing
      fairness definitions (equalized odds / demographic parity, fair regression,
      AP-based audits).
\item Specified confounders (object scale; occlusion where annotated) and the
      ICON$^2$ + Mann--Whitney statistical protocol.
\item Added a discussion of the practical safety implications of $1$--$3$\,pp
      IoU gaps for depth / TTC estimation.
\item Clarified that race/skin tone is \emph{not} one of the two audited axes,
      with justification and future-work plans.
\end{itemize}

%=====================================================================
\section*{Response to Reviewer 1 (Score: 2 --- Accept)}
%=====================================================================

\subsection*{R1.1 --- Theoretical relation of Localization Parity to existing fairness definitions}

\textbf{Comment.}
\textit{Clarify how Localization Parity relates theoretically to existing
fairness definitions for regression or structured prediction; discuss
mathematical properties.}

\textbf{Response.}
We agree this strengthens the contribution. Section~3.2 of the camera-ready
paper now situates LP explicitly:
\begin{itemize}[leftmargin=1.2em,itemsep=2pt]
\item Relative to \textbf{demographic parity / equalized odds}, LP replaces a
      binary detection outcome with a continuous quality score (IoU), evaluated
      only on the detected set---analogous to equalized opportunity applied to
      localization quality.
\item Relative to \textbf{fair regression}, LP instantiates bounded disparity
      in expected loss with $\ell = 1 - \mathrm{IoU}$ on matched boxes.
\item Relative to \textbf{AP-based audits}, LP isolates the localization
      component that AP conflates with classification confidence and ranking.
\end{itemize}
We also state clearly that LP is a group-level criterion (not individual
fairness or calibration), chosen because the safety harm of interest is
group-structured. Orthogonality to detection parity is formalized in
Eq.~(1) and the surrounding text.

\subsection*{R1.2 --- Confounders and statistical methodology}

\textbf{Comment.}
\textit{Specify the confounders considered and the statistical methodology
used.}

\textbf{Response.}
Section~4.4 now states this explicitly:
\begin{itemize}[leftmargin=1.2em,itemsep=2pt]
\item \textbf{Primary confounder:} object scale (ground-truth box height),
      because female and child pedestrians tend to occupy smaller image regions.
\item \textbf{Secondary confounder:} occlusion, used where annotations permit.
\item \textbf{Method:} ICON$^2$ stratified binning into short / medium / tall
      height terciles within each dataset; per-group mean IoU recomputed inside
      each stratum.
\item \textbf{Test:} two-sided Mann--Whitney $U$ (non-parametric; IoU
      distributions are non-Gaussian), significance at $p<0.05$.
\item \textbf{Result:} gender gap survives in $35/36$ stratified cells; age gap
      in $24/30$.
\end{itemize}

\subsection*{R1.3 --- Practical safety implications of $1$--$3$\,pp IoU gaps}

\textbf{Comment.}
\textit{Discuss how localization disparities of approximately $1$--$3$
percentage points propagate into downstream perception or planning errors.}

\textbf{Response.}
New Section~3.3 provides a first-order geometric argument: under pinhole
projection, relative depth error scales with relative box-height error. For a
pedestrian at $30$\,m with an $80$\,px box height, a $\sim$2\,pp mean IoU
deficit (roughly $2$--$4$\,px systematic height / center error) corresponds to
approximately $0.7$--$1.5$\,m of depth bias---enough to alter time-to-collision
estimates at urban speeds. Because trajectory and collision modules are
typically nonlinear in box geometry, group-level mean gaps of this magnitude
remain safety-relevant even when every detection clears the $0.5$ IoU threshold.
We openly note that a full closed-loop propagation study is future work, now
listed in the Conclusion.

\subsection*{R1.4 --- Generalization beyond gender/age and autonomous driving}

\textbf{Comment.}
\textit{Unclear whether the framework generalizes to other protected attributes
or domains beyond autonomous driving.}

\textbf{Response.}
LP is domain-agnostic: any detector that outputs boxes and any protected
attribute with instance-level labels can be audited under Eq.~(1). Empirically,
this paper studies gender and age in driving scenes because those labels are
available at scale. The Ethics and Conclusion sections now state clearly that
skin tone, intersectional groups, disability, and non-driving domains
(surveillance, robotics) are planned extensions, and explain why skin tone was
excluded from the present audit (insufficient labeled sample / coarse
categories in our pilot).

\vspace{0.4em}
\noindent We thank Reviewer~1 for the positive overall assessment and for
highlighting the value of transparently reporting the negative generalization
result for DWBBR; that discussion is retained and sharpened in Section~5.2.

%=====================================================================
\section*{Response to Reviewer 2 (Score: 1 --- Weak Accept)}
%=====================================================================

\subsection*{R2.1 --- Datasets must be specified}

\textbf{Comment.}
\textit{No dataset is mentioned; data source, size, and composition must be
specified.}

\textbf{Response.}
Addressed throughout. The abstract and Section~4.1 now name all four datasets
with approximate sizes and composition:
\begin{itemize}[leftmargin=1.2em,itemsep=2pt]
\item \textbf{BDD100K} --- US dashcam; $\sim$100{,}000 pedestrian boxes; binary
      gender labels.
\item \textbf{CityPersons} --- $\sim$19{,}000 boxes / 2{,}975 images.
\item \textbf{EuroCity-Day} and \textbf{EuroCity-Night} --- subsets of EuroCity
      Persons ($\sim$238{,}200 boxes / 47{,}300 images / 31 cities).
\end{itemize}
Phase~2--3 training uses a 1{,}693-image BDD100K subset (423-image validation).

\subsection*{R2.2 --- Detector families must be named}

\textbf{Comment.}
\textit{Detector families are not named in the abstract.}

\textbf{Response.}
Addressed. The abstract and Section~4.2 now name:
\begin{itemize}[leftmargin=1.2em,itemsep=2pt]
\item \textbf{YOLOX} (single-stage, anchor-free),
\item \textbf{RetinaNet} (single-stage, anchor-based),
\item \textbf{Cascade R-CNN} (two-stage),
\end{itemize}
with \textbf{YOLOv8n} as the trainable backbone for the DWBBR intervention.

\subsection*{R2.3 --- Demographic axes / race--skin tone}

\textbf{Comment.}
\textit{Unclear whether race/skin tone was one of the two demographic axes;
axes should be stated explicitly.}

\textbf{Response.}
The two axes are \textbf{gender} (female vs.\ male) and \textbf{age} (child
vs.\ adult). Race / skin tone is \textbf{not} one of them. This is now stated
explicitly in the abstract, Introduction (contribution~2), Section~4.3, and
Ethics. We briefly explain the exclusion (pilot sample too small / categories
too coarse) and list skin-tone LP as future work.

\subsection*{R2.4 --- Conference template}

\textbf{Comment.}
\textit{The PanAfriCon conference template has not been used; please reformat
to the required Springer LaTeX template.}

\textbf{Response.}
Completed. The camera-ready manuscript is prepared with the official Springer
LNCS class (\texttt{llncs.cls}, bibliography style \texttt{splncs04}). Source
and compiled PDF are included in the submission zip.

%=====================================================================
\section*{Closing}
%=====================================================================

We believe the revised manuscript fully addresses the reviewers' concerns while
preserving the core contributions that both reviews found valuable: a new
fairness criterion, a broad confounder-controlled audit, and transparent
reporting of a mitigation that succeeds in training but fails to generalize.
We are grateful for the opportunity to improve the paper and look forward to
presenting this work at PanAfriCon AI 2026.

\vspace{1.2em}
\noindent Best regards,\\[0.6em]
Abel Adissu, Anwar Gashaw Yimam, Dawit Kindea, Beakal Gizachew, and Elefelious Getachew\\
Addis Ababa University\\
\texttt{abel.adissu-ug@aau.edu.et},
\texttt{anwar.gashaw-ug@aau.edu.et},\\
\texttt{dawit.kindea-ug@aau.edu.et},
\texttt{beakal.gizachew@aau.edu.et},\\
\texttt{elefelious.getachew@aau.edu.et}

\end{document}
